\documentclass[a4paper,final]{article}
\usepackage[pagewise]{lineno}
\usepackage{amssymb}
\usepackage{enumerate}
\usepackage{amsmath}
\usepackage{amsthm}
%\usepackage{a4wide}
\usepackage{setspace}
\usepackage{makeidx}
\usepackage[skip=5pt plus1pt]{parskip}
\usepackage{natbib}
\usepackage{latexsym}
\usepackage[dvipsnames]{xcolor}
\usepackage{tabularx}
\usepackage{enumitem}
\usepackage{makecell}

\usepackage[breaklinks,colorlinks,linkcolor=black,citecolor=black,urlcolor=gray,hypertexnames=false]{hyperref} % the addition of hypertexnames is for autonum
\usepackage[capitalize]{cleveref}
\usepackage{tikz}
\usepackage{xfrac}
\usepackage{wrapfig}
\usepackage{mathpazo}
\usepackage{lscape}
\usepackage{stmaryrd}
%\usepackage{caption}\captionsetup{font=footnotesize} %% CCM: it was breaking labels that come after table!
%\usepackage[USenglish]{babel}
\usepackage{graphicx}

\usepackage{showlabels}


\usepackage{multirow}


% Define "struts" as suggested by Claudio Beccari in
% a piece in TeX and TUG News, Vol. 2, 1993.
\newcommand\Tstrut{\rule{0pt}{2.6ex}}       % "top" strut
\newcommand\Bstrut{\rule[-0.9ex]{0pt}{0pt}} % "bottom" strut
\newcommand{\TBstrut}{\Tstrut\Bstrut} % top&bottom struts


\newtheorem{theorem}{Theorem}[section]
\newtheorem{axiom}{Axiom}
\newtheorem{metatheorem}[theorem]{Metatheorem}
\newtheorem{definition}[theorem]{Definition}
\newtheorem{metadefinition}[theorem]{Metadefinition}
\newtheorem{proposition}[theorem]{Proposition}
\newtheorem{metaproposition}[theorem]{Metaproposition}
\newtheorem{lemma}[theorem]{Lemma}
\newtheorem{metalemma}[theorem]{Metalemma}
\newtheorem{corollary}[theorem]{Corollary}
\newtheorem{metacorollary}[theorem]{Metacorollary}
\newtheorem{example}[theorem]{Example}
\newtheorem{ruleinf}{Rule of Inference}
\newtheorem{sublemma}{Sublemma}[theorem]
%\doublespacing

%\linenumbersx

\DeclareMathOperator*{\argmin}{\arg\,\min}
\DeclareMathOperator*{\argmax}{\arg\,\max}

\DeclareMathOperator*{\arginf}{\arg\,\inf}
\DeclareMathOperator*{\argsup}{\arg\,\sup}
\DeclareMathOperator{\posi}{posi}

% Load amsfonts for \mathbb (optional, but recommended)
%\usepackage{amsfonts}
% Load mathbbol to extend \mathbb to Greek letters




	\newcommand{\quotient}[2]{\left.#1\middle/#2\right.}

\newcommand\F{\mathcal{F}}
\newcommand\X{\mathcal{X}}
\newcommand\Y{\mathcal{Y}}
\newcommand\A{\mathcal{A}}
\renewcommand\P{\mathbb{P}} % impreicse prob 
\newcommand\cl{\mathrm{cl}}
%\newcommand\PowN{\imprecpickstrat} % power set of set of prob strategies
\newcommand\Uset{\mathbb{U}} % set of utility functions
\newcommand\Prob{\mathcal{P}} % set of all probs

\newcommand\C{\mathcal{C}}
\newcommand\Exp{\mathsf{Exp}}
\newcommand\RExp{\mathrm{R}\Exp} % Risk-weighted expected utility 
\newcommand\EU{\mathrm{EU}}
\newcommand\MEU{\mathrm{MEU}}
\newcommand\REU{\mathrm{REU}}
\newcommand\EAd{\mathrm{EAd}}
\renewcommand\O{\mathcal{O}}
\newcommand\U{\mathfrak{U}} % utility random variable 
\newcommand\Uwald{\mathcal{U}} % utility random variable 
\newcommand\ut{u}
\newcommand\Maxi{\Gamma\mathrm{M}}
\newcommand\Maximality{\mathrm{Max}}
\newcommand\Maximin{\Gamma}
\newcommand\GD{\mathrm{GD}}

%\newcommand\D{\mathsf{DecProblems}}
\newcommand{\D}{\mathcal{D}}
	\newcommand{\Decs}{\mathcal{D}}
		\newcommand{\mass}{\mathcal{M}}
	
\renewcommand\S{\mathcal{S}}
\newcommand\s{\mathsf{s}}
\renewcommand\c{\mathsf{c}} % choice function

\newcommand{\n}{\mathsf{n}}
\renewcommand{\nu}{\n}
\newcommand\Nu{\mathcal{N}}
\newcommand{\imprecpickstrat}{\mathbb{N}}

\newcommand{\IB}{\mathbb{B}}
\newcommand{\ID}{\mathbb{M}}
\newcommand{\IP}{\P}
\newcommand{\pb}{b}
%\newcommand{b}{}
%\renewcommand{\color}[2]{}
\newcommand{\Dom}{\mathsf{Dom}}


\renewcommand{\Re}{\mathbb{R}}


\renewcommand{\mathsterling}{\text{\normalfont\textsterling}} % to be able to use \pounds in math mode
\newcommand{\million}{\mathrm{m}}   % million shorthand


% CCM added
\usepackage[textsize=footnotesize,obeyFinal,backgroundcolor=orange!50,bordercolor=gray!20,linecolor=gray!20]{todonotes} % use "\PassOptionsToPackage{disable}{todonotes}" to disable in a document
\newcommand{\todoinfo}[2][]{\todo[inline,{#1}]{#2}}
\newcommand{\todoold}[2][]{\todo[backgroundcolor=white,bordercolor=orange!10,linecolor=gray!10, #1,caption={},textcolor=gray]{Pre-rev: #2}}
\newcommand{\comment}[2][]{\todoold[#1]{#2}}
\newcommand{\todooldinfo}[2][]{\todoold[#1]{#2}}
\usepackage{booktabs}  % For improved table formatting


\makeatletter
\@ifclasswith{article}{final}{
	% Redefine commands or settings for the final version.
\renewcommand{\color}[1]{}
}{
	% Alternative definition for when final is not present.
%	\renewcommand{\mycommand}{Draft Version}
}
\makeatother


\usepackage{mathtools} % using it for definition of \Set
\usepackage{autonum} % only number equations referenced to
\usepackage{bm} % get bold greek symbols
\newenvironment{colored}[1]{\leavevmode\color{#1}}{}
\usepackage{nicefrac}

\newcommand{\Strategies}{\S}
\newcommand{\Conv}{\mathsf{ConvHull}}

\newcommand\SetDelimiter[1][]{
	\nonscript\,#1\vert \allowbreak \nonscript\,\mathopen{}}
\providecommand\given{\SetDelimiter}
\DeclarePairedDelimiterX\Set[1]{\lbrace}{\rbrace}%
{ \renewcommand\given{\SetDelimiter[\delimsize]} #1 } % Spacing around \given in a set. \Set{x\given blah} has good spacing. Can also use starred version \Set*{x\given blah} when blah is multilined so it'll stretch to contain everything. 
\DeclarePairedDelimiterX\seq[1]{\langle}{\rangle}%
{ \renewcommand\given{\SetDelimiter[\delimsize]} #1 } % Spacing around \given in a set. \Set{x\given blah} has good spacing. Can also use starred version \Set*{x\given blah} when blah is multilined so it'll stretch to contain everything. 
\DeclarePairedDelimiter\abs{\lvert}{\rvert}%

%\newrobustcmd*{\citefirstlastauthor}{\AtNextCite{\DeclareNameAlias{labelname}{given-family}}\citeauthor}
\newcommand*\diff{\mathop{}\!\mathrm{d}}
%\newcommand{\vec}{def}

\renewcommand{\emptyset}{\varnothing}
\renewcommand{\leq}{\leqslant}
\renewcommand{\geq}{\geqslant}




\newenvironment{CCM rewritten}
{\begingroup\color{blue}} % Begin environment
{\endgroup}              % End environment


\usepackage{versions}
\excludeversion{infversion}

\begin{document}

\title{Positive-Density Self-Undermining of Risk-Weighted Expected Utility in Finite Dimensions}
\author{}
\date{}
\maketitle

\section{Proof A (using Branden)}

\begin{theorem}[Positive-density self-undermining of risk-weighted expected utility]
\label{thm:reu-positive-density}

Let $\Omega$ be a finite state space, and let $p$ be a probability
function on $\Omega$ such that
\[
p(\omega)>0
\]
for every $\omega\in\Omega$. Fix $l<h$, and identify an act with its
utility profile
\[
u\in[l,h]^\Omega.
\]

Let
\[
\Decs:=([l,h]^\Omega)^2
\]
be the space of ordered binary decision problems. Thus, if
\[
D=(u,v)\in\Decs,
\]
the acts available in $D$ are $u$ and $v$. A picking strategy is a
measurable function
\[
\s:\Decs\longrightarrow[l,h]^\Omega
\]
such that
\[
\s(u,v)\in\{u,v\}
\]
for every $(u,v)\in\Decs$.

Let
\[
r:[0,1]\longrightarrow[0,1]
\]
be continuous and strictly increasing, with
\[
r(0)=0,
\qquad
r(1)=1,
\qquad
r\in C^2((0,1)),
\qquad
r'(x)>0
\]
for every $x\in(0,1)$.

For an act $u$, define
\[
\REU_{r,p}(u)
:=
l+\int_l^h
r\bigl(p(\{\omega\in\Omega:u(\omega)>x\})\bigr)
\,\mathrm{d}x.
\]

Let $\mu$ be a Borel probability measure on $\Decs$. Suppose that
$\mu$ is absolutely continuous with respect to Lebesgue measure on
$\Decs$, and that its density is positive almost everywhere on
\[
((l,h)^\Omega)^2.
\]

Assume complete independence between the state and the decision
problem, so that the joint probability on $\Omega\times\Decs$ is
\[
b=p\times\mu.
\]
For a measurable picking strategy $\s$, define
\[
f_\s(x)
:=
\int_{\Decs}
p\bigl(\{\omega\in\Omega:\s(D)(\omega)>x\}\bigr)
\,\mu(\mathrm{d}D),
\]
and define its higher-order risk-weighted expected utility by
\[
\REU_{r,b}(\s)
:=
l+\int_l^h r(f_\s(x))\,\mathrm{d}x.
\]

Suppose there is a non-empty proper event
\[
A\subsetneq\Omega
\]
such that
\[
r\bigl(p(A)\bigr)\neq p(A).
\]

Then every measurable picking strategy $\s$ satisfying
\[
\s(D)\in
\arg\max_{a\in D}\REU_{r,p}(a)
\]
for every $D\in\Decs$ is strictly suboptimal at the higher-order
level. More precisely, there is a measurable picking strategy $\s'$
such that
\[
\REU_{r,b}(\s')>\REU_{r,b}(\s).
\]

Moreover, $\s'$ can be chosen so that
\[
\mu\left(
\left\{
D\in\Decs:
\s'(D)\notin
\arg\max_{a\in D}\REU_{r,p}(a)
\right\}
\right)>0.
\]
\end{theorem}


\begin{proof}
Write
\[
\mathcal D:=([l,h]^\Omega)^2.
\]
For an act \(u\in[l,h]^\Omega\), define its tail function by
\[
H_u(x)
:=
p\bigl(\{\omega\in\Omega:u(\omega)>x\}\bigr)
=
\sum_{\omega\in\Omega}
p(\omega)\mathbf 1_{\{u(\omega)>x\}}.
\]
For a measurable picking strategy \(s\), define
\[
f_s(x)
:=
\int_{\mathcal D}H_{s(D)}(x)\,\mu(\mathrm dD).
\]
Then, by complete independence,
\[
\REU_{r,p\times\mu}(s)
=
l+\int_l^h r(f_s(x))\,\mathrm dx.
\]

\medskip
\noindent
\textbf{Step 1: Regularity of the induced tail function.}

For each \(\omega\in\Omega\), the selected utility
\[
D\longmapsto s(D)(\omega)
\]
has an atomless distribution under \(\mu\). Indeed, if
\(N\subseteq[l,h]\) has Lebesgue measure zero, then
\[
\{(u,v):s(u,v)(\omega)\in N\}
\subseteq
\{(u,v):u(\omega)\in N\}
\cup
\{(u,v):v(\omega)\in N\}.
\]
Both sets on the right have \(\mu\)-measure zero because \(\mu\) is
absolutely continuous with respect to Lebesgue measure. It follows that
\[
x\longmapsto
\mu\{D:s(D)(\omega)>x\}
\]
is continuous for every \(\omega\), and hence \(f_s\) is continuous.

Moreover,
\[
0<f_s(x)<1
\]
for every \(x\in(l,h)\). To see that \(f_s(x)>0\), consider the set
\[
G_x^+
:=
\left\{
(u,v)\in((l,h)^\Omega)^2:
u(\omega)>x
\text{ and }
v(\omega)>x
\text{ for every }\omega\in\Omega
\right\}.
\]
This is a non-empty open subset of the interior of \(\mathcal D\). On
\(G_x^+\), whichever act \(s\) selects has utility greater than \(x\)
in every state, and hence
\[
H_{s(D)}(x)=1.
\]
Since the density of \(\mu\) is positive almost everywhere,
\[
\mu(G_x^+)>0,
\]
and therefore \(f_s(x)>0\).

Similarly, the set
\[
G_x^-
:=
\left\{
(u,v)\in((l,h)^\Omega)^2:
u(\omega)<x
\text{ and }
v(\omega)<x
\text{ for every }\omega\in\Omega
\right\}
\]
is non-empty and open. On this set,
\[
H_{s(D)}(x)=0.
\]
Thus \(\mu(G_x^-)>0\), and consequently \(f_s(x)<1\).

\medskip
\noindent
\textbf{Step 2: The shadow value.}

Fix \(y_0\in(l,h)\), and define
\[
W_s(y)
:=
\int_{y_0}^{y}r'(f_s(x))\,\mathrm dx.
\]
Then
\[
W_s'(y)=r'(f_s(y))>0.
\]

Define the shadow value of an act \(u\) by
\[
L_s(u)
:=
\sum_{\omega\in\Omega}
p(\omega)W_s(u(\omega)).
\]

The significance of \(L_s\) is that it gives the first-order effect on
the higher-order value of replacing one selected act by another. To
see this, note that
\begin{align}
\int_l^h r'(f_s(x))H_u(x)\,\mathrm dx
&=
\sum_{\omega\in\Omega}
p(\omega)
\int_l^{u(\omega)}
r'(f_s(x))\,\mathrm dx
\\
&=
L_s(u)+C_s,
\end{align}
where
\[
C_s
:=
\sum_{\omega\in\Omega}
p(\omega)
\int_l^{y_0}r'(f_s(x))\,\mathrm dx
\]
does not depend on \(u\). Therefore, for any two acts \(u\) and \(v\),
\[
\int_l^h
r'(f_s(x))
\bigl(H_v(x)-H_u(x)\bigr)
\,\mathrm dx
=
L_s(v)-L_s(u).
\]

\medskip
\noindent
\textbf{Step 3: A strict reversal between first-order REU and the
shadow value.}

By assumption, there is a non-empty proper event
\[
A\subsetneq\Omega
\]
such that
\[
r(p(A))\neq p(A).
\]
Put
\[
P:=p(A).
\]
Since \(p(\omega)>0\) for every \(\omega\), we have
\[
0<P<1.
\]

Consider the two-parameter family of acts
\[
u^{x,y}(\omega)
:=
\begin{cases}
x, & \omega\in A,\\
y, & \omega\notin A,
\end{cases}
\qquad
l<x,y<h,
\qquad
x>y.
\]
For such an act,
\[
H_{u^{x,y}}(z)
=
\begin{cases}
1, & z<y,\\
P, & y\leq z<x,\\
0, & z\geq x.
\end{cases}
\]
Hence
\[
\REU_{r,p}(u^{x,y})
=
r(P)x+\bigl(1-r(P)\bigr)y.
\]
The common lower-bound term cancels in this expression.

Meanwhile,
\[
L_s(u^{x,y})
=
P W_s(x)+(1-P)W_s(y).
\]

We claim that there are \(x>y\) and \(x'>y'\) such that
\[
\REU_{r,p}(u^{x,y})
>
\REU_{r,p}(u^{x',y'})
\]
but
\[
L_s(u^{x,y})
<
L_s(u^{x',y'}).
\]

Suppose not. Then the gradients of these two functions of \((x,y)\)
must be positively collinear at every point of the open region
\[
\{(x,y)\in(l,h)^2:x>y\}.
\]
Indeed, if the gradients were not positively collinear at some point,
there would be a direction \(d\in\mathbb R^2\) that increased the first
function and decreased the second. A sufficiently small movement in
that direction would produce the required strict reversal.

The two gradients are
\[
\bigl(r(P),1-r(P)\bigr)
\]
and
\[
\bigl(PW_s'(x),(1-P)W_s'(y)\bigr).
\]
Positive collinearity therefore requires
\[
\frac{W_s'(x)}{W_s'(y)}
=
\frac{r(P)(1-P)}
     {(1-r(P))P}
=:C
\]
for every \(x>y\).

Choose
\[
x>y>z.
\]
Then
\[
\frac{W_s'(x)}{W_s'(y)}=C
\qquad\text{and}\qquad
\frac{W_s'(y)}{W_s'(z)}=C,
\]
so
\[
\frac{W_s'(x)}{W_s'(z)}=C^2.
\]
But applying the same relation directly to \(x>z\) gives
\[
\frac{W_s'(x)}{W_s'(z)}=C.
\]
Therefore
\[
C^2=C.
\]
Since \(W_s'>0\), we have \(C>0\), and hence
\[
C=1.
\]
It follows that
\[
r(P)(1-P)
=
(1-r(P))P,
\]
and therefore
\[
r(P)=P.
\]
This contradicts
\[
r(p(A))\neq p(A).
\]

Thus there are two acts \(a,c\in((l,h)^\Omega)\) such that
\[
\REU_{r,p}(a)>\REU_{r,p}(c)
\]
but
\[
L_s(c)>L_s(a).
\]

\medskip
\noindent
\textbf{Step 4: Thickening the reversal to an open set of decision
problems.}

The two acts \(a\) and \(c\) obtained above are constant on \(A\) and
on its complement. Although such acts lie in a lower-dimensional
subset of the full act space, the two inequalities are strict.
Risk-weighted expected utility is continuous in the utility profile,
and \(L_s\) is continuous because \(W_s\) is continuously
differentiable.

Therefore there are open neighbourhoods
\[
O_a,O_c\subseteq(l,h)^\Omega
\]
of \(a\) and \(c\), respectively, and some \(\eta>0\), such that for
every \(u\in O_a\) and \(v\in O_c\),
\[
\REU_{r,p}(u)>\REU_{r,p}(v)
\]
and
\[
L_s(v)-L_s(u)\geq\eta.
\]

We may shrink \(O_a\) and \(O_c\) so that their closures are compact
subsets of \((l,h)^\Omega\). Thus there is a compact interval
\[
J\subseteq(l,h)
\]
containing every coordinate of every act in
\[
O_a\cup O_c.
\]

Put
\[
N:=O_a\times O_c.
\]
Then \(N\) is a non-empty open subset of the full-dimensional
decision space
\[
((l,h)^\Omega)^2.
\]
Because the density of \(\mu\) is positive almost everywhere,
\[
\mu(N)>0.
\]

For every decision problem
\[
D=(u,v)\in N,
\]
we have
\[
\REU_{r,p}(u)>\REU_{r,p}(v).
\]
Since \(s\) selects an \(\REU_{r,p}\)-maximising act,
\[
s(D)=u
\]
throughout \(N\). However,
\[
L_s(v)-L_s(s(D))
=
L_s(v)-L_s(u)
\geq\eta.
\]

\medskip
\noindent
\textbf{Step 5: A needle variation.}

Since \(\mu\) is absolutely continuous, it is nonatomic. We may
therefore choose a measurable set
\[
E\subseteq N
\]
with arbitrarily small positive measure
\[
\mu(E)=\varepsilon.
\]

Define a new measurable picking strategy \(t\) by
\[
t(D)
:=
\begin{cases}
v, & D=(u,v)\in E,\\
s(D), & D\notin E.
\end{cases}
\]
Thus \(t\) differs from \(s\) only on \(E\), and on \(E\) it selects
the act with the larger shadow value.

Write
\[
g_E(x):=f_t(x)-f_s(x).
\]
Since \(s\) and \(t\) differ only on a set of measure \(\varepsilon\),
and since every tail function takes values in \([0,1]\),
\[
|g_E(x)|\leq\varepsilon
\]
for every \(x\).

Moreover,
\[
g_E(x)=0
\]
outside \(J\). Below the lower endpoint of \(J\), both selected acts
have tail value \(1\); above the upper endpoint of \(J\), both have
tail value \(0\).

Define
\[
B_r(y,z)
:=
r(y)-r(z)-r'(z)(y-z).
\]
Then
\[
r(f_t(x))-r(f_s(x))
=
r'(f_s(x))g_E(x)
+
B_r(f_t(x),f_s(x)).
\]
Therefore
\begin{align}
\REU_{r,p\times\mu}(t)
-
\REU_{r,p\times\mu}(s)
&=
\int_l^h r'(f_s(x))g_E(x)\,\mathrm dx
\\
&\quad+
\int_l^h
B_r(f_t(x),f_s(x))
\,\mathrm dx.
\label{eq:reu-needle-decomposition}
\end{align}

By the definition of \(f_s\) and \(f_t\), and by Fubini's theorem,
\begin{align}
\int_l^h r'(f_s(x))g_E(x)\,\mathrm dx
&=
\int_E
\int_l^h
r'(f_s(x))
\bigl(H_{t(D)}(x)-H_{s(D)}(x)\bigr)
\,\mathrm dx
\,\mu(\mathrm dD)
\\
&=
\int_E
\bigl(L_s(t(D))-L_s(s(D))\bigr)
\,\mu(\mathrm dD).
\end{align}
On \(E\),
\[
L_s(t(D))-L_s(s(D))\geq\eta.
\]
Hence
\[
\int_l^h r'(f_s(x))g_E(x)\,\mathrm dx
\geq
\eta\varepsilon.
\]

It remains to bound the second term in
\eqref{eq:reu-needle-decomposition}. Since \(f_s\) is continuous and
\[
0<f_s(x)<1
\]
for every \(x\in(l,h)\), its image on the compact interval \(J\) lies
in a compact subinterval of \((0,1)\). For sufficiently small
\(\varepsilon\), the same is true of \(f_t\), because
\[
|f_t(x)-f_s(x)|\leq\varepsilon.
\]

Since \(r\in C^2((0,1))\), there is some \(M<\infty\) such that
\[
|r''(z)|\leq M
\]
for every value of \(z\) lying between \(f_s(x)\) and \(f_t(x)\), with
\(x\in J\). Taylor's theorem therefore gives
\[
|B_r(f_t(x),f_s(x))|
\leq
\frac{M}{2}|g_E(x)|^2
\leq
\frac{M}{2}\varepsilon^2.
\]
Because \(g_E\) vanishes outside \(J\),
\[
\left|
\int_l^h
B_r(f_t(x),f_s(x))
\,\mathrm dx
\right|
\leq
\frac{M}{2}\lambda(J)\varepsilon^2,
\]
where \(\lambda(J)\) is the length of \(J\).

Thus, for some constant \(C_0<\infty\),
\[
\REU_{r,p\times\mu}(t)
-
\REU_{r,p\times\mu}(s)
\geq
\eta\varepsilon-C_0\varepsilon^2.
\]
For sufficiently small \(\varepsilon>0\),
\[
\eta\varepsilon-C_0\varepsilon^2>0.
\]
Therefore
\[
\REU_{r,p\times\mu}(t)
>
\REU_{r,p\times\mu}(s).
\]

Finally,
\[
\mu(E)=\varepsilon>0,
\]
and on every \(D\in E\), the strategy \(t\) selects the act \(v\)
even though
\[
\REU_{r,p}(u)>\REU_{r,p}(v).
\]
Thus \(t\) fails to select an \(\REU_{r,p}\)-maximising act on a set
of positive \(\mu\)-measure, as required.
\end{proof}
\section{Alternative proof}
\begin{theorem}[Positive-density self-undermining of risk-averse REU]
\label{thm:reu-positive-density}
Let $\Omega$ be a finite state space with at least two elements, and let
$p$ be a probability function on $\Omega$ such that
\[
p(\omega)>0
\]
for every $\omega\in\Omega$. Let
\[
\A:=[-1,1]^\Omega,
\qquad
\X:=\A^2,
\]
where $D=(a,b)\in\X$ is an ordered binary decision problem. A
measurable picking strategy is a measurable function $\s:\X\to\A$
such that
\[
\s(a,b)\in\{a,b\}
\]
for every $(a,b)\in\X$.

Let $r:[0,1]\to[0,1]$ be continuous and strictly increasing, with
\[
r(0)=0,
\qquad
r(1)=1,
\qquad
r\in C^2((0,1)),
\qquad
r'(q)>0
\]
for every $q\in(0,1)$. Suppose moreover that
\[
r(q)<q
\]
for every $q\in(0,1)$.

Let $\mu$ be a Borel probability measure on $\X$ that is absolutely
continuous with respect to Lebesgue measure and whose density is
positive almost everywhere on
\[
\bigl((-1,1)^\Omega\bigr)^2.
\]
For $a\in\A$ and $x\in[-1,1]$, write
\[
p(a>x):=p\bigl(\{\omega\in\Omega:a(\omega)>x\}\bigr),
\]
and define
\[
\REU_{r,p}(a)
:=-1+\int_{-1}^{1}r\bigl(p(a>x)\bigr)\diff x.
\]
Under complete independence, define the risk-weighted expected utility
of a picking strategy $\s$ by
\[
\REU_{r,p\times\mu}(\s)
:=-1+\int_{-1}^{1}
r\left(
\int_\X p\bigl(\s(D)>x\bigr)\,\mu(\diff D)
\right)\diff x.
\]

If $\s$ selects an $\REU_{r,p}$-maximising act from every decision
problem, then there is a measurable picking strategy $\s'$ such that
\[
\REU_{r,p\times\mu}(\s')
>
\REU_{r,p\times\mu}(\s).
\]
Moreover,
\[
\mu\bigl(\{D\in\X:\s'(D)\neq\s(D)\}\bigr)>0,
\]
and, on this positive-measure set, $\s'$ selects the unique
non-$\REU_{r,p}$-maximising act.
\end{theorem}

\begin{proof}
We identify each number in $[-1,1]$ with the corresponding constant
act. Write
\[
\Phi(a):=\REU_{r,p}(a).
\]
For the given strategy $\s$, define
\[
Q(x):=
\int_\X p\bigl(\s(D)>x\bigr)\,\mu(\diff D),
\qquad x\in[-1,1].
\]

We first show that $Q$ is continuous. For each $\omega\in\Omega$, let
\[
Y_\omega(D):=\s(D)(\omega).
\]
For every $t\in[-1,1]$,
\[
\{D:Y_\omega(D)=t\}
\subseteq
\{(a,b):a(\omega)=t\}
\cup
\{(a,b):b(\omega)=t\}.
\]
The two sets on the right are Lebesgue-null hyperplanes, and therefore
have $\mu$-measure zero. Thus $Y_\omega$ has an atomless distribution
under $\mu$. Since
\[
Q(x)=\sum_{\omega\in\Omega}
p(\omega)\,\mu\bigl(Y_\omega>x\bigr),
\]
it follows that $Q$ is continuous.

Let
\begin{align}
B_+
&:=\{(a,b)\in\X:0<a(\omega)<1\text{ and }0<b(\omega)<1
\text{ for every }\omega\in\Omega\},\\
B_-
&:=\{(a,b)\in\X:-1<a(\omega)<0\text{ and }-1<b(\omega)<0
\text{ for every }\omega\in\Omega\}.
\end{align}
Both are non-empty open subsets of the interior of $\X$, so
\[
\mu(B_+)>0
\qquad\text{and}\qquad
\mu(B_-)>0.
\]
On $B_+$, whichever act $\s$ selects is strictly positive in every
state; on $B_-$, whichever act it selects is strictly negative in every
state. Hence
\[
0<Q(0)<1.
\]
By continuity, there are $\eta,\kappa>0$ such that
\[
I:=[-\eta,\eta]\subset(-1,1)
\]
and
\[
\kappa\leq Q(x)\leq1-\kappa
\qquad
\text{for every }x\in I.
\]
Define
\[
w(x):=r'(Q(x)),
\qquad x\in I.
\]
Then $w$ is continuous and strictly positive on $I$.

Fix a non-empty proper event $A\subsetneq\Omega$, and put
\[
\alpha:=p(A).
\]
Since $p$ has full support and $r(\alpha)<\alpha$, we may choose
$c$ such that
\[
r(\alpha)<c<\alpha.
\]
For $\varepsilon>0$, define the act
\[
g_\varepsilon(\omega)
:=\varepsilon\bigl(\mathbf 1_A(\omega)-c\bigr).
\]
Thus
\[
g_\varepsilon(\omega)
=
\begin{cases}
\varepsilon(1-c),&\omega\in A,\\
-\varepsilon c,&\omega\notin A.
\end{cases}
\]
For all sufficiently small $\varepsilon>0$, both the constant-zero act
and $g_\varepsilon$ lie in the interior of $\A$, and all values taken
by these acts lie in the interior of $I$. Moreover,
\[
\Phi(0)=0,
\]
while
\begin{align}
\Phi(g_\varepsilon)
&=-1
+\int_{-1}^{-\varepsilon c}1\diff x
+\int_{-\varepsilon c}^{\varepsilon(1-c)}
r(\alpha)\diff x\\
&=\varepsilon\bigl(r(\alpha)-c\bigr)<0.
\end{align}
Thus, in the decision problem $(0,g_\varepsilon)$, the constant-zero
act is the unique $\REU_{r,p}$-maximiser.

We next establish the inequality that will make it beneficial to replace
the constant-zero act by nearby gamble-like acts on a small open set of
decision problems. Put
\[
J_\varepsilon
:=\int_I
w(x)\bigl(p(g_\varepsilon>x)-p(0>x)\bigr)\diff x.
\]
Apart from endpoints,
\[
p(g_\varepsilon>x)-p(0>x)
=
\begin{cases}
-(1-\alpha),&-\varepsilon c<x<0,\\
\alpha,&0<x<\varepsilon(1-c),\\
0,&\text{otherwise}.
\end{cases}
\]
Consequently,
\[
J_\varepsilon
=-(1-\alpha)\int_{-\varepsilon c}^{0}w(x)\diff x
+\alpha\int_0^{\varepsilon(1-c)}w(x)\diff x.
\]
By continuity of $w$ at zero,
\begin{align}
\lim_{\varepsilon\downarrow0}\frac{J_\varepsilon}{\varepsilon}
&=w(0)\bigl(-(1-\alpha)c+\alpha(1-c)\bigr)\\
&=w(0)(\alpha-c)>0.
\end{align}
Fix $\varepsilon>0$ sufficiently small that
\[
J_\varepsilon>0
\]
and all values taken by $0$ and $g_\varepsilon$ lie in the interior of
$I$. Set
\[
D^*:=\{0,g_\varepsilon\}.
\]

The map $\Phi:\A\to\Re$ is continuous. Indeed, if $a_n\to a$, then
\[
r\bigl(p(a_n>x)\bigr)\longrightarrow r\bigl(p(a>x)\bigr)
\]
for every $x$ except possibly the finitely many values
\[
\{a(\omega):\omega\in\Omega\},
\]
and dominated convergence applies.

For acts $a,b$ whose values lie in the interior of $I$, define
\[
G(a,b)
:=\int_I
w(x)\bigl(p(b>x)-p(a>x)\bigr)\diff x.
\]
This map is continuous. To see this, let
\[
W(y):=\int_0^y w(x)\diff x,
\qquad y\in I.
\]
Then
\[
G(a,b)
=
\sum_{\omega\in\Omega}
p(\omega)\bigl(W(b(\omega))-W(a(\omega))\bigr).
\]
At $D^*$,
\[
\Phi(0)>\Phi(g_\varepsilon)
\qquad\text{and}\qquad
G(0,g_\varepsilon)=J_\varepsilon>0.
\]
Therefore there are an open neighbourhood $V$ of $D^*$ and a number
$\gamma>0$ such that, for every $D=(a,b)\in V$,
\[
\Phi(a)>\Phi(b)
\qquad\text{and}\qquad
G(a,b)>\gamma.
\]
We may take $V$ small enough that every value of every act appearing
in $V$ lies in the interior of $I$.

Since $D^*$ lies in the interior of $\X$, sufficiently small open balls
about $D^*$ are contained in $V$. Every such ball has positive
$\mu$-measure, because the density of $\mu$ is positive almost
everywhere on the interior. Their measures converge to zero as their
radii converge to zero, because $\mu$ is absolutely continuous and
therefore has no atoms.

Let
\[
M:=
\sup_{q\in[\kappa/2,\,1-\kappa/2]}
|r''(q)|<\infty.
\]
Choose an open ball $\O$ about $D^*$, with $\O\subseteq V$, such that,
writing
\[
m:=\mu(\O),
\]
we have
\[
0<m<\frac{\kappa}{2}
\qquad\text{and}\qquad
M\eta m<\gamma.
\]
For $D=\{a_D,b_D\}\in\O$, the first act $a_D$ is the unique
$\REU_{r,p}$-maximiser, and hence
\[
\s(D)=a_D.
\]
Define
\[
\s'(D)
:=
\begin{cases}
b_D,&D=(a_D,b_D)\in\O,\\
\s(D),&D\notin\O.
\end{cases}
\]
This is a measurable picking strategy. It takes the second act in every
decision problem in $\O$.

Define
\[
\delta(x)
:=
\int_\O
\bigl(p(b_D>x)-p(a_D>x)\bigr)\,\mu(\diff D).
\]
The aggregate survival function generated by $\s'$ is
\[
Q_{\s'}(x)=Q(x)+\delta(x).
\]
Because all values of the acts in $\O$ lie in $I$,
\[
\delta(x)=0
\qquad\text{for }x\notin I,
\]
and, for every $x$,
\[
|\delta(x)|\leq m.
\]
By Fubini's theorem,
\begin{align}
\int_I r'(Q(x))\delta(x)\diff x
&=\int_\O G(a_D,b_D)\,\mu(\diff D)\\
&\geq\gamma m.
\end{align}

For $x\in I$, both $Q(x)$ and $Q(x)+\delta(x)$ lie in
\[
[\kappa/2,1-\kappa/2].
\]
Taylor's theorem therefore gives
\[
r(Q(x)+\delta(x))-r(Q(x))
\geq
r'(Q(x))\delta(x)-\frac{M}{2}\delta(x)^2.
\]
Since $\delta$ vanishes outside $I$,
\begin{align}
&\REU_{r,p\times\mu}(\s')
-\REU_{r,p\times\mu}(\s)\\
&\quad=
\int_I
\bigl(r(Q(x)+\delta(x))-r(Q(x))\bigr)\diff x\\
&\quad\geq
\gamma m-\frac{M}{2}\int_I\delta(x)^2\diff x\\
&\quad\geq
\gamma m-M\eta m^2\\
&\quad>0.
\end{align}
Thus
\[
\REU_{r,p\times\mu}(\s')
>
\REU_{r,p\times\mu}(\s).
\]

Finally, $\s'$ differs from $\s$ at every decision problem in $\O$, and
only there. Hence
\[
\mu\bigl(\{D:\s'(D)\neq\s(D)\}\bigr)
=\mu(\O)=m>0.
\]
On $\O$, $\s'$ selects $b_D$, which is the unique
non-$\REU_{r,p}$-maximising act. This proves the result.
\end{proof}


\end{document}
